\section*{Bab \ref{ch:g9:solids} --- Bangun Ruang dan Volume}

\begin{solution}{exo:g9:solids:1}
\omterm{def:g5:solids:def}{Balok}: $V = 4 \times 5 \times 12 = 240$.

\omterm{def:g7:areas:prism}{Tabung}: $V = \pi r^2 h = \pi \times 9 \times 7 = 63\pi \approx 198$.
\end{solution}

\begin{solution}{exo:g9:solids:2}
\omterm{def:g8:solids:def}{Kerucut}: $V = \frac13 \pi r^2 h = \frac13 \pi \times 25 \times 12 =
100\pi \approx 314$.

\omterm{def:g8:solids:def}{Limas}: luas alasnya $B = 8 \times 3 = 24$, jadi
$V = \frac13 \times 24 \times 10 = 80$.
\end{solution}

\begin{solution}{exo:g9:solids:3}
\omterm{def:g6:measure:volume}{Volumenya}: $V = \frac43 \pi \times 6^3 = \frac43 \pi \times 216 = 288\pi$.

Luas permukaannya: $4\pi \times 6^2 = 144\pi$.
\end{solution}

\begin{solution}{exo:g9:solids:4}
\omterm{def:g3:shapes:shapes}{Jari-jari} \omterm{prop:g9:solids:sections}{penampangnya}: $\sqrt{r^2 - d^2} = \sqrt{100 - 64} =
\sqrt{36} = 6$.
\end{solution}

\begin{solution}{exo:g9:solids:5}
\omterm{def:g9:thales:scaling}{Faktor skala} dari seluruh \omterm{def:g8:solids:def}{kerucutnya} ke yang kecil:
$k = \frac{9}{12} = \frac34$. \omterm{def:g3:shapes:shapes}{Jari-jari} \omterm{prop:g9:solids:sections}{penampangnya}:
$8 \times \frac34 = 6$.

\omterm{def:g6:measure:volume}{Volume} seluruh \omterm{def:g8:solids:def}{kerucutnya}: $\frac13\pi \times 64 \times 12 = 256\pi$.
\omterm{def:g6:measure:volume}{Volume} \omterm{def:g8:solids:def}{kerucut} kecilnya: $256\pi \times \left(\frac34\right)^3 =
256\pi \times \frac{27}{64} = 108\pi \approx 339$.
\end{solution}

\begin{solution}{exo:g9:solids:6}
\omterm{def:g6:measure:volume}{Volume} bolanya: $\frac43\pi \times 2^3 = \frac{32\pi}{3}$ cm$^3$.
\omterm{def:g6:measure:volume}{Volume} ini menyebar pada luas alas $\pi \times 3^2 = 9\pi$ cm$^2$,
jadi permukaannya naik sebesar
\[
\frac{32\pi/3}{9\pi} = \frac{32}{27} \approx 1.2 \text{ cm}
\]
(kira-kira $12$ mm).
\end{solution}

\begin{solution}{exo:g9:solids:7}
Cetakan kecilnya adalah pengecilan cetakan besarnya dengan
$k = \frac36 = \frac12$, jadi \omterm{def:g6:measure:volume}{volumenya}
$\left(\frac12\right)^3 = \frac18$ \omterm{def:g6:measure:volume}{volume} yang besar: jadi resepnya
mengisi $8$ cetakan kecil.
\end{solution}

\begin{solution}{exo:g9:solids:8}
$V = \frac13 \times 230^2 \times 147 = \frac13 \times 52\,900 \times 147
= 52\,900 \times 49 = 2\,592\,100$ m$^3$ --- kira-kira $2.6$ juta meter
kubik.
\end{solution}

\begin{solution}{exo:g9:solids:9}
\emph{1.} Airnya membentuk \omterm{def:g8:solids:def}{kerucut} yang terskala $k = \frac12$
(puncaknya di bawah, setengah tingginya), jadi \omterm{def:g6:measure:volume}{volumenya}
$\left(\frac12\right)^3 = \frac18$ \omterm{def:g6:measure:volume}{volume} corongnya: yaitu
seperdelapan, jauh lebih sedikit daripada setengahnya!

\emph{2.} Air sampai ketinggian $d$ membentuk \omterm{def:g8:solids:def}{kerucut} yang terskala
$k = \frac{d}{12}$, dengan \omterm{def:g6:measure:volume}{volume} $k^3$ kali \omterm{def:g6:measure:volume}{volume} corongnya.
Setengah \omterm{def:g6:measure:volume}{volumenya} berarti $k^3 = \frac12$, yaitu
$\left(\frac{d}{12}\right)^3 = \frac12$. Lalu
$\frac{d}{12} = \sqrt[3]{0.5} \approx 0.7937$, jadi
$d \approx 9.5$ cm: separuh kedua \omterm{def:g6:measure:volume}{volumenya} hanya menempati
$2.5$ cm teratas --- bagian yang lebar pada \omterm{def:g8:solids:def}{kerucutnya} memuat
sebagian besar airnya.
\end{solution}

\begin{solution}{pb:g9:solids:1}
\textbf{1.} $V_{\text{tabung}} = \pi r^2 \times 2r = 2\pi r^3$.

\textbf{2.} $\dfrac{V_{\text{bola}}}{V_{\text{tabung}}}
= \dfrac{\frac43 \pi r^3}{2 \pi r^3} = \dfrac23$. Faktor $\pi$ dan
$r^3$ saling meniadakan: perbandingannya \emph{semesta} --- benar
untuk sebutir kelereng dan untuk sebuah planet. Yaitu hubungan antara
bentuk, bukan antara bilangan: persis jenis kebenaran yang layak
dipahat pada batu.

\textbf{3.} Bola: $4\pi r^2$. \omterm{def:g7:areas:prism}{Tabung}: selimutnya
$2\pi r \times 2r = 4\pi r^2$, ditambah dua tutupnya
$2 \times \pi r^2$: seluruhnya $6\pi r^2$. Perbandingannya:
$\frac{4\pi r^2}{6\pi r^2} = \frac23$ --- yaitu dua pertiga yang sama,
untuk permukaan. (Dan sebagai bonus: luas bolanya persis sama dengan
luas selimut \omterm{def:g7:areas:prism}{tabungnya}.)

\textbf{4.} Sisanya: $2\pi r^3 - \frac43\pi r^3 =
\frac23 \pi r^3$. Dua \omterm{def:g8:solids:def}{kerucut} berjari-jari $r$ dan tinggi $r$:
$2 \times \frac13 \pi r^2 \times r = \frac23 \pi r^3$. Sama.

\textbf{5.} Bolanya: $\frac43 \pi \times 12^3 \approx 7\,200$
cm$^3$; kotaknya: $2\pi \times 12^3 \approx 10\,900$ cm$^3$. Yang
kosong: sepertiga kotaknya, kira-kira $33\,\%$ --- yang dijamin
pertanyaan 2, berapa pun besar bolanya.

\textbf{6.} Setengah bola:
$\frac12 \times \frac43 \pi \times 10^3 = \frac23 \pi \times
1\,000 \approx 2\,094$ cm$^3 \approx 2.1$ L.

\textbf{7.} \omterm{def:g8:solids:def}{Kerucut}: $\frac13 \pi \times 10^2 \times 10 =
\frac{1\,000\pi}{3} \approx 1\,047$ cm$^3$. Setengah bola:
$\frac{2\,000\pi}{3} \approx 2\,094$ cm$^3$. \omterm{def:g7:areas:prism}{Tabung}:
$1\,000\pi \approx 3\,142$ cm$^3$. Perbandingannya:
$\frac{1000\pi}{3} : \frac{2000\pi}{3} : \frac{3000\pi}{3}
= 1 : 2 : 3$ persis.

\textbf{8.} $\frac43 \pi \times 27 = \pi \times 9 \times h$
memberi $h = \frac43 \times 3 = 4$ cm persis.

\textbf{9.} \omterm{def:g3:shapes:shapes}{Jari-jarinya}: $\frac{6\,371}{1\,737} \approx 3.67$.
\omterm{def:g6:measure:volume}{Volume} menskala seperti kubik (\cref{thm:g9:solids:scaling}):
$3.67^3 \approx 49$. Kira-kira lima puluh Bulan muat di dalam Bumi.

\textbf{10.} Permukaan menskala seperti kuadrat: $3.67^2 \approx
13.5$. Melipatduakan \omterm{def:g3:shapes:shapes}{jari-jari} balon mengalikan \omterm{def:g6:measure:volume}{volumenya} dengan
$2^3 = 8$ tetapi permukaannya hanya dengan $2^2 = 4$: jadi ketika
benda membesar, \omterm{def:g6:measure:volume}{volume} (berat, isi, panas yang dihasilkan) melampaui
permukaan (kulit, bahan, panas yang hilang) --- itulah hukum
kuadrat--kubik.

\textbf{11.} Beratnya: $\times 10^3 = 1\,000$. \omterm{prop:g9:solids:sections}{Penampang}
tulangnya: $\times 10^2 = 100$. Tekanannya --- yaitu berat tiap
luas tulangnya: $\times \frac{1000}{100} = 10$. Tulang yang dibangun
untuk tekanan manusia menerima sepuluh kali lebih besar: kerangka
raksasanya patah di bawah beratnya sendiri. Raksasa mustahil secara
geometris; hewan yang besar perlu tulang yang tebalnya tidak
\omterm{def:g9:linfunc:linear}{sebanding} (bandingkan kaki gajah dengan kaki kijang).

\textbf{12.} Beratnya: $\div 100^3 = 10^6$. Kekuatannya (yaitu
\omterm{prop:g9:solids:sections}{penampang} otot): $\div 100^2 = 10^4$. Kekuatan terhadap beratnya:
$\times \frac{10^6}{10^4} = 100$. Semut yang mengangkat lima puluh
kali beratnya bukan atlet super --- ia hanya kecil; seorang manusia
yang dikerutkan sampai seukuran semut dapat berbuat serupa.

\textbf{13.} $\dfrac{S}{V} = \dfrac{4\pi r^2}{\frac43\pi r^3}
= \dfrac{3}{r}$: untuk $r = 1$ perbandingannya $3$, untuk $r = 2$
perbandingannya $1.5$ --- membagi dua ukurannya melipatduakan
permukaan yang tersedia tiap satuan \omterm{def:g6:measure:volume}{volume} penghasil panas. Tubuh
yang kecil membuang panas: tikusnya harus makan terus-menerus,
beruangnya dapat berhibernasi, dan bayinya --- bola kecil dengan
$\frac Sr$ yang besar --- perlu topinya.

\textbf{14.} $\left(\frac54\right)^3 = \frac{125}{64} \approx
1.95$: jeruk yang besar memuat hampir \emph{dua kali} buahnya untuk
harga yang sama. Belilah \omterm{def:g3:shapes:shapes}{jari-jarinya}.

\textbf{15.} Pahatannya berkata: bola $= \frac23$ \omterm{def:g7:areas:prism}{tabungnya}
menurut \omterm{def:g6:measure:volume}{volume}, \emph{dan} $\frac23$ \omterm{def:g7:areas:prism}{tabungnya} menurut permukaan
seluruhnya --- yaitu dua perbandingan yang persis dan semesta yang
menghubungkan \omterm{def:g5:solids:def}{bangun ruang} paling bundar dengan yang paling
sederhana. Jawaban untuk Cicero: karena sebuah teorema adalah
satu-satunya monumen yang tidak dikikis tentara maupun abad ---
mesin Archimedes ikut terbakar di kota Syracusa, tetapi dua
pertiganya masih persis dua pertiga.

\end{solution}
